\chapter{State, Unitaries and Measurement}\label{ch:state}

\begin{etymology}
\textbf{State} comes from the Latin \emph{status}, from \emph{stare}, ``to
stand''. The condition something stands in at a given moment, whatever path
brought it there.

\textbf{Unitary} comes from \emph{unitas}, ``oneness'', from \emph{unus},
``one''. Something that stays whole. A unitary transformation can change the form
of a state but keeps it whole.

\textbf{Measurement} comes from \emph{mensura}, from \emph{metiri}, ``to measure'',
to judge by comparison with a standard.
\end{etymology}

Chapter~\ref{ch:qubit} used the words state, unitary and measurement loosely.
This chapter pins them down for a single qubit:

\begin{quote}
A unitary operation changes the state without destroying quantum information: it
moves amplitudes and phases around but keeps the length of the state vector. A
measurement does the opposite: it returns one result and replaces the
superposition with it.
\end{quote}

\section{The state as a vector}

Draw a vector of length 1 in the plane, at 45 degrees. Its shadows on the two
axes have length $\cos 45^\circ = \sin 45^\circ = 1/\sqrt{2} \approx 0.707$. Tilt it
to 30 degrees and the shadows become $0.866$ and $0.5$. Square them and you get
$0.75$ and $0.25$, which add to 1, because the vector still has length 1.

That is already the picture of a qubit. The vector is a combination of the two
axes, and tilting it moves weight from one to the other while keeping the total
fixed. The squared shadows are the probabilities of the two outcomes.

What the plane cannot show is phase. The state of a qubit is a vector in
$\mathbb{C}^2$: two complex numbers $(\alpha, \beta)$, each with a size and an
angle. Sizes become probabilities; the angle between them, the relative phase,
controls how the components combine. A real number describes a line. A pair of
reals describes a plane. A single complex number already describes a plane, and
two of them describe a qubit.

\section{The Bloch sphere}

Two complex numbers are four real numbers. Normalization,
$|\alpha|^2 + |\beta|^2 = 1$, removes one. The global phase, which no measurement
can see (Chapter~\ref{ch:hilbert}), removes another. Two remain, and two
numbers describe a point on a sphere. Every state of one qubit can be written as
\[
  \ket{\psi} = \cos\tfrac{\theta}{2}\,\ket{0} + e^{i\varphi}\sin\tfrac{\theta}{2}\,\ket{1},
\]
where $\theta$ is the latitude, measured from the north pole, and $\varphi$ the
longitude.

\begin{figure}[htbp]
\centering
\begin{tikzpicture}[scale=2.1]
  \draw[black!50, line width=0.7pt] (0,0) circle (1);
  \draw[black!35, dashed] (1,0) arc (0:180:1 and 0.3);
  \draw[black!50] (-1,0) arc (180:360:1 and 0.3);
  \draw[black!35, -{Stealth[length=1.6mm]}] (0,0) -- (0,1.25) node[above, font=\sffamily\small, text=black] {$\ket{0}$};
  \draw[black!35] (0,0) -- (0,-1) node[below, font=\sffamily\small, text=black] {$\ket{1}$};
  \draw[black!35, -{Stealth[length=1.6mm]}] (0,0) -- (-0.55,-0.42) node[below left, font=\sffamily\scriptsize, text=black!60] {$x$};
  \draw[black!35, -{Stealth[length=1.6mm]}] (0,0) -- (1.25,0) node[right, font=\sffamily\scriptsize, text=black!60] {$y$};
  % latitude of the state theta = 60 deg: z = cos 60 = 0.5, radius sin 60 = 0.866
  \draw[hfHistory, line width=0.7pt] (0.866,0.5) arc (0:-180:0.866 and 0.26);
  \draw[hfHistory, dashed, line width=0.5pt] (0.866,0.5) arc (0:180:0.866 and 0.26);
  \draw[hfConcept, line width=1.2pt, -{Stealth[length=2.2mm]}] (0,0) -- (-0.70,0.35);
  \fill[hfConcept] (-0.70,0.35) circle (0.9pt);
  \node[font=\sffamily\small, text=hfConcept] at (-0.88,0.22) {$\ket{\psi}$};
  \draw[black!60] (0,0.32) arc (90:150:0.32 and 0.3);
  \node[font=\sffamily\scriptsize] at (-0.1,0.42) {$\theta$};
  \node[hf/note, text=hfHistory, anchor=west] at (1.05,0.55) {same $|\alpha|^2$, any $\varphi$};
\end{tikzpicture}
\caption{The Bloch sphere. $\ket{0}$ is at the north pole and $\ket{1}$ at the
south pole. The latitude fixes the probabilities; moving around the circle of
latitude changes only the relative phase.}
\end{figure}

\begin{history}{1946 \textperiodcentered{} Stanford}
Felix Bloch introduced the sphere to describe nuclear spins in his work on
nuclear magnetic induction, the basis of NMR and, much later, of MRI. Henri
Poincaré had drawn the same sphere in 1892 for the polarization of light. The
mathematics is identical, because both are two-level systems.
\end{history}

Fix $|\alpha|^2 = 0.75$ and $|\beta|^2 = 0.25$ and you fix the latitude: the
state is at $\theta = 60^\circ$. But that does not fix the state. The relative
phase $\varphi$ moves it around that circle of latitude without touching the
probabilities. The state
$\sqrt{0.75}\,\ket{0} + 0.5\,\ket{1}$, with $\varphi = 0$, sits at the point
$(0.866, 0, 0.5)$. The sphere invents nothing; it only draws the two things a
qubit state carries, how much and which phase, in a form the eye can follow.

\section{Matrices}

If a state is a vector, an operation on a state is something that takes a vector
and returns another. In linear algebra that object is a \emph{matrix}. It can
rotate, stretch, flip or mix the components of a vector, depending on its
entries. A matrix is not just a table of numbers; it is the algebraic form of a
transformation (Chapter~\ref{ch:linear}).

\section{Unitary transformations}

The state has length 1, and that is not a technicality: it is what makes the
probabilities add to 100\%. So a physical operation cannot change a state any way
it likes. If it stretched or shrank the vector, the probabilities would stop
adding up. The operations allowed are the ones that preserve length.

In the real plane, the transformations that change direction but not length are
rotations. In complex space, they are called \emph{unitary}, and they satisfy
\[
  U^\dagger U = I ,
\]
where $U^\dagger$ is the conjugate transpose of $U$ and $I$ is the identity. A
state evolves as $\ket{\psi'} = U\ket{\psi}$, and the condition guarantees that
$\braket{\psi'|\psi'} = \braket{\psi|U^\dagger U|\psi} = \braket{\psi|\psi} = 1$.
It also guarantees that $U$ has an inverse, $U^\dagger$, so every unitary step can
be undone.

\begin{concept}{Quantum gate}
A unitary transformation used as an elementary step of a circuit, usually acting
on one or two qubits. Deutsch introduced quantum circuits made of gates in 1989.

\textbf{What it fixes:} complicated evolutions can be built from a small menu of
simple, calibrated operations. A few gates (for example $H$, $T$ and CNOT) are
enough to approximate any unitary.

\textbf{What it costs:} every gate on real hardware has an error, and a useful
algorithm needs millions or billions of them.
\end{concept}

The gates of Chapter~\ref{ch:qubit}, in matrix form:
\[
  H = \frac{1}{\sqrt 2}\begin{pmatrix} 1 & 1 \\ 1 & -1 \end{pmatrix},
  \qquad
  Z = \begin{pmatrix} 1 & 0 \\ 0 & -1 \end{pmatrix},
  \qquad
  X = \begin{pmatrix} 0 & 1 \\ 1 & 0 \end{pmatrix}.
\]
$X$ swaps $\ket{0}$ and $\ket{1}$, the quantum NOT. $Z$ flips the sign of
$\ket{1}$. On the Bloch sphere, each one is a rotation. And $HZH = X$: the
interference trick of Chapter~\ref{ch:qubit}, sign flip between two
Hadamards, is a NOT gate in disguise.

\section{Measurement}

Unitary evolution is smooth, reversible and deterministic. Measurement is none of
those things. Measuring a qubit in the basis $\{\ket{0}, \ket{1}\}$ returns 0
with probability $|\alpha|^2$ and 1 with probability $|\beta|^2$ (Born's rule,
Chapter~\ref{ch:quantum}). Afterwards, the state is no longer
$\alpha\ket{0} + \beta\ket{1}$: it is $\ket{0}$ or $\ket{1}$, whichever came out.
Von Neumann wrote this \emph{projection postulate} into his 1932 book. Measuring
again gives the same answer; the amplitudes are gone.

So you never see a superposition. You see one outcome per run, and to estimate the
amplitudes you have to prepare the same state and measure it many times. The goal
of an algorithm is not to look at the superposition but to arrange the unitaries
so that, at the moment of measurement, the useful answer has a high probability.

Measuring in a different basis is the same as rotating first and then measuring in
the standard one. Apply $H$ and measure, and you are asking whether the qubit was
$\ket{+} = (\ket{0} + \ket{1})/\sqrt2$ or $\ket{-} = (\ket{0} - \ket{1})/\sqrt2$,
a question about the relative phase that the standard measurement cannot answer.
Every measurement answers some questions and destroys the answers to others.

\section{Why complex amplitudes?}

The interference of Chapter~\ref{ch:qubit} only used signs, $+$ and $-$. Real
amplitudes with signs already reinforce and cancel. So why complex?

Signs give two phases, $0$ and $180$ degrees. Complex numbers give a continuous
circle of them, and quantum mechanics uses all of it: the time evolution
$e^{-i\hat Ht}$ turns phases continuously, and the longitude $\varphi$ of the Bloch
sphere would collapse to a single great circle without it.

For a long time physicists suspected that the complex numbers might be a
convenience, and that an equivalent theory with real numbers only could exist.
For single systems it can. In 2021, Marc-Olivier Renou and colleagues showed that
in networks with independent sources, real-number quantum theory makes different
predictions from the standard one, and experiments in 2022 sided with complex
numbers. The $i$ in quantum mechanics is not bookkeeping.
